\documentclass[12pt]{article}
\usepackage[utf8]{inputenc}
\usepackage[T1]{fontenc}
\usepackage{hyperref}
\title{Pana Language Resources for K-12 Curriculum Development}
\author{Compiled by Claude Sonnet 5.5 (claude-sonnet-5-5) for Andr\'as Kornai}
\date{}
\begin{document}
\maketitle
\section{Introduction}
Pana (ISO 639-3 pnq; also called Sama) is a Gur (Niger-Congo) language of northwest Burkina Faso and Mali, along the Sourou River. Speaker counts differ by source, from about 5,000 to about 8,000. The language is classed as threatened by the Endangered Languages Project.
\section{Primary Educational Resources}
None found. No primers or school programmes appeared in search results.
\section{Linguistic and Reference Materials}
\begin{itemize}
\item Klaus Beyer (2006), \emph{La langue pana (Burkina Faso et Mali): description linguistique, lexique, textes}, Gur Monographs 7, R\"udiger K\"oppe Verlag. It includes a grammar sketch, a Pana--French lexicon and a French--Pana list, and three texts. It is in French and is sold in print. \url{https://www.koeppe.de/titel_la-langue-pana-burkina-faso-et-mali}
\item Wikipedia, Pana language (Gur). \url{https://en.wikipedia.org/wiki/Pana_language_(Gur)}
\item Endangered Languages Project entry, Pana (Burkina Faso). \url{https://endangeredlanguages.com/lang/4825?hl=fr}
\item Endangered Languages Project bibliography for Pana. \url{https://www.endangeredlanguages.com/lang/4825/bibliography}
\end{itemize}
\section{Religious and Community Resources}
\begin{itemize}
\item Joshua Project, Pana in Mali (demographic profile). \url{https://joshuaproject.net/index.php/people_groups/18893/ML}
\end{itemize}
\section{Digital and Audio Resources}
None found.
\section{Curriculum Development Context}
The only substantial resource found is Beyer's French-language description and lexicon, which could seed an orthography and word list. No existing mother-tongue teaching materials were found.
\end{document}